
        \documentclass[fleqn]{article}      
        \usepackage{amsmath}
        \usepackage{fontspec}
        \usepackage{kotex} % 한국어 지원  

        \begin{document}      
        % Easy Problems
\noindent 1. 다음 행렬의 합을 구하시오:  
\[
A = \begin{pmatrix}
1 & 2 \\
3 & 4
\end{pmatrix}, \quad 
B = \begin{pmatrix}
5 & 6 \\
7 & 8
\end{pmatrix}
\]  
\\
A + B = ? \\\\

\noindent 2. 행렬 \( C \)의 두배를 구하시오:  
\[
C = \begin{pmatrix}
2 & 4 \\
6 & 8
\end{pmatrix}
\]  
\\
2C = ? \\\\

\noindent 3. 행렬 \( A \)와 \( B \)의 곱을 구하시오:  
\[
A = \begin{pmatrix}
1 & 0 \\
0 & 1
\end{pmatrix}, \quad 
B = \begin{pmatrix}
2 & 3 \\
4 & 5
\end{pmatrix}
\]  
\\
AB = ? \\\\

\noindent 4. 다음 행렬의 행렬식 값은?  
\[
D = \begin{pmatrix}
1 & 2 \\
3 & 4
\end{pmatrix}
\]  
\\
\(|D| = ?\) \\\\

\noindent 5. 다음 행렬의 전치행렬을 구하시오:  
\[
E = \begin{pmatrix}
1 & 2 & 3 \\
4 & 5 & 6
\end{pmatrix}
\]  
\\
E^T = ? \\\\

% Medium Problems
\noindent 6. 두 행렬의 덧셈을 구하시오:  
\[
F = \begin{pmatrix}
1 & 3 & 5 \\
2 & 4 & 6 \\
\end{pmatrix}, \quad 
G = \begin{pmatrix}
7 & 8 & 9 \\
10 & 11 & 12 
\end{pmatrix}
\]  
\\
F + G = ? \\\\

\noindent 7. 다음 행렬의 스칼라 곱을 구하시오:  
\[
H = \begin{pmatrix}
2 & -1 \\
0 & 3
\end{pmatrix}
\]  
\\
3H = ? \\\\

\noindent 8. 행렬 \( I \)의 전치행렬을 구하시오:  
\[
I = \begin{pmatrix}
1 & 2 \\
3 & 4 \\
5 & 6
\end{pmatrix}
\]  
\\
I^T = ? \\\\

\noindent 9. 다음 행렬의 행렬식을 구하시오:  
\[
J = \begin{pmatrix}
2 & 3 \\
1 & 4
\end{pmatrix}
\]  
\\
\(|J| = ?\) \\\\

\noindent 10. 행렬 \( K \)와 \( L \)의 곱을 구하시오:  
\[
K = \begin{pmatrix}
1 & 2 \\
3 & 4
\end{pmatrix}, \quad 
L = \begin{pmatrix}
5 & 6 \\
7 & 8
\end{pmatrix}
\]  
\\
KL = ? \\\\

% Difficult Problems
\noindent 11. 행렬 \( M \)의 고유값을 구하시오:  
\[
M = \begin{pmatrix}
3 & 1 \\
0 & 2
\end{pmatrix}
\]  
\\
고유값 = ? \\\\

\noindent 12. 다음 행렬의 역행렬을 구하시오:  
\[
N = \begin{pmatrix}
4 & 7 \\
2 & 6
\end{pmatrix}
\]  
\\
\( N^{-1} = ? \) \\\\

\noindent 13. 다음 행렬의 고유벡터를 구하시오:  
\[
O = \begin{pmatrix}
2 & 1 \\
1 & 2
\end{pmatrix}
\]  
\\
고유벡터 = ? \\\\

\noindent 14. 두 행렬 \( P \)와 \( Q \)의 행렬식의 곱을 구하시오:  
\[
P = \begin{pmatrix}
1 & 2 \\
3 & 4
\end{pmatrix}, \quad 
Q = \begin{pmatrix}
2 & 0 \\
1 & 3
\end{pmatrix}
\]  
\\
\(|P||Q| = ?\) \\\\

\noindent 15. 다음 연립방정식을 행렬을 이용하여 풀어보시오:  
\[
\begin{cases}
2x + 3y = 5 \\
4x + 9y = 13
\end{cases}
\]  
\\
\((x, y) = ?\) \\\\\

% Answers
1. \(\begin{pmatrix} 6 & 8 \\ 10 & 12 \end{pmatrix}\) \\\\
2. \(\begin{pmatrix} 4 & 8 \\ 12 & 16 \end{pmatrix}\) \\\\
3. \(\begin{pmatrix} 2 & 3 \\ 4 & 5 \end{pmatrix}\) \\\\
4. \(-2\) \\\\
5. \(\begin{pmatrix} 1 & 4 \\ 2 & 5 \\ 3 & 6 \end{pmatrix}\) \\\\
6. \(\begin{pmatrix} 8 & 11 & 14 \\ 12 & 15 & 18 \end{pmatrix}\) \\\\
7. \(\begin{pmatrix} 6 & -3 \\ 0 & 9 \end{pmatrix}\) \\\\
8. \(\begin{pmatrix} 1 & 3 & 5 \\ 2 & 4 & 6 \end{pmatrix}\) \\\\
9. \(5\) \\\\
10. \(\begin{pmatrix} 19 & 22 \\ 43 & 50 \end{pmatrix}\) \\\\
11. \( \{3, 2\} \) \\\\
12. \(\begin{pmatrix} 3/10 & -7/10 \\ -1/5 & 2/5 \end{pmatrix}\) \\\\
13. \( \begin{pmatrix} 1 \\ -1 \end{pmatrix} \) \\\\
14. \( -6 \) \\\\
15. \((1, 1)\) \\\\\      
        \end{document} 
    